\chapter{Trading sur les nouvelles et les événements}\label{ch:hf:news-and-event-trading}

À huit heures et demie, le premier vendredi du mois, un chiffre est publié. En moins d'une milliseconde, les contrats à terme s'échangent à Chicago, et les teneurs de marché
qui ne voulaient pas en faire partie retirent leurs cotations une seconde plus tôt. Dans le modèle de publication de ce chapitre, le niveau le plus rapide des traders, situé
à 50~microsecondes du chiffre, capte 88~\% de tout ce que les cotations périmées de la publication abandonnent, et le niveau suivant se partage le reste~; un teneur de marché qui
laisse ses cotations en place perd 58 lots-tick par publication, et celui qui les retire puis les remet cinq secondes plus tard gagne 29~\% de plus sur les deux minutes suivantes que
celui qui n'est jamais parti.

\section{Chiffres programmés et course à la publication}

Les publications économiques (livre~2, chapitre~31) arrivent à des heures connues, avec une prévision consensuelle~; le prix bouge avec la
surprise statistique. Leurs éditeurs contrôlent qui voit le chiffre en premier.

\begin{definition}[Lock-up de publication]\label{def:hf:news-and-event-trading:lockup}
Un \emph{lock-up de publication}\index{lock-up de publication} (release lock-up) est un dispositif par lequel une agence statistique ou un éditeur transmet le chiffre à des journalistes ou à des fournisseurs de
données avant sa publication officielle, sous embargo et sans communication extérieure, de sorte que les dépêches et les flux soient prêts
à l'heure de la publication.
\end{definition}

\begin{definition}[Course à la publication]\label{def:hf:news-and-event-trading:race}
Une \emph{course à la publication}\index{course à la publication} (release race) est la \omterm{def:hf:latency-arbitrage-and-its-defence:arb}{course à la latence} (chapitre~9) déclenchée par une publication programmée~: les traders qui reçoivent et analysent le
chiffre en premier prennent les cotations que ce chiffre a rendues périmées, dans chaque instrument qui en dépend.
\end{definition}

L'issue de la course dépend de qui reçoit le chiffre en premier, et de la vitesse à laquelle les cotations qu'il rend périmées sont retirées. Hu, Pan et Wang ont étudié une
période allant de 2007 à juin 2013, pendant laquelle certains traders à haute vitesse recevaient l'indice de confiance des consommateurs du Michigan deux secondes avant sa publication officielle~:
les transactions se concentraient fortement sur ces deux secondes, et la formation des prix prenait moins de 200 millisecondes. En comparant avec d'autres publications et
avec la période postérieure à la fin de l'accord, ils ont conclu que la diffusion échelonnée pouvait avoir réduit, plutôt qu'accru, l'avantage informationnel des traders
les plus rapides, et rendu la formation des prix plus efficiente.

\begin{dated}{2026-09}{Accès anticipé, faux titre et tweet}\label{dat:hf:news-and-event-trading:events}
Le 8 juillet 2013, le procureur général de l'État de New York a annoncé que Thomson Reuters avait accepté de mettre fin immédiatement à la pratique consistant à fournir
à des traders à haute fréquence certains résultats d'enquête sur la consommation, susceptibles de faire bouger les marchés, deux secondes avant ses autres abonnés. Le 23 avril 2013, peu après
13~h, un faux tweet publié depuis le compte piraté de l'Associated Press a signalé des explosions à la Maison-Blanche~; le Dow Jones Industrial Average, le
Nasdaq et le S\&P 500 ont plongé pendant quelques minutes, jusqu'à ce que les traders constatent que l'information était fausse. Le 29 septembre 2018, la SEC a annoncé qu'Elon Musk
et Tesla verseraient chacun une amende de 20~millions de dollars et que Musk quitterait la présidence du conseil, en règlement des poursuites engagées au sujet de ses tweets, dans lesquels il affirmait pouvoir retirer Tesla
de la cote à 420~dollars l'action, les financements étant assurés.
\end{dated}

\section{Titres lisibles par machine}

Les dépêches lisibles par machine (livre~8, chapitre~17) transmettent titres et chiffres dans des champs qu'un programme peut analyser sans les lire. Le trader qui
analyse en premier gagne la course~; le trader qui analyse mal la perd. Une stratégie qui capte 2~ticks lorsqu'elle lit correctement un titre et en perd 10
lorsqu'elle le lit mal gagne 1.76~tick par titre avec un taux d'erreur de 2~\%, et perd au-delà de 20~\%~: une erreur d'analyse, et non la vitesse, est souvent
la contrainte déterminante dès qu'une entreprise est assez rapide.

\section{Réseaux sociaux et événements non programmés}

Les événements non programmés surviennent sans prévision consensuelle, souvent par des canaux conçus pour des personnes plutôt que pour des machines, et parfois de manière fausse. Le
piratage d'avril 2013 présenté dans l'encadré daté a fait bouger l'ensemble du marché américain pendant quelques minutes~; un tweet du directeur général d'une entreprise a fait bouger son cours de bourse et s'est soldé par un
règlement. Un teneur de marché ne peut pas gagner une course dont il ignore qu'elle a commencé~; ce qu'il peut faire, c'est reconnaître qu'elle a commencé, et protéger ses cotations.

\section{Le protocole d'événement du teneur de marché}

\begin{definition}[Protocole d'événement]\label{def:hf:news-and-event-trading:protocol}
Un \emph{protocole d'événement}\index{protocole d'événement} (event protocol) est le calendrier écrit d'un teneur de marché pour un événement connu~: quand annuler ou élargir ses cotations avant lui,
quelles données doivent être mises à jour avant la reprise de la cotation, quand et avec quelle largeur réinscrire ses ordres, et quelles limites s'appliquent jusqu'au retour à une cotation normale.
\end{definition}

Le modèle du chapitre (\cref{lst:hf:news-and-event-trading:race})~: la surprise du chiffre déplace la fair value du contrat à terme de trois ticks par
écart type~; dix niveaux de cinquante lots chacun, fournis par d'autres teneurs qui mettent à jour leurs cotations 20~millisecondes après la publication, auxquels s'ajoutent les vingt lots de notre teneur de marché à chacun des
trois premiers niveaux~; quatre niveaux de preneurs de liquidité, à 0.05, 1, 5 et 50~millisecondes, de deux cents lots chacun. Un lot exécuté à $k$ ticks de l'ancien prix
rapporte $J-k$ ticks pour un mouvement de $J$.

\begin{center}
\small
\begin{tabular}{@{}lrrrr@{}}
\toprule
latence des preneurs & 0.05~ms & 1~ms & 5~ms & 50~ms\\
\midrule
part de la valeur des cotations périmées captée & 88.1~\% & 11.6~\% & 0.3~\% & 0\\
\bottomrule
\end{tabular}
\end{center}

La course est à rafle presque totale~: 223 lots-tick de valeur périmée par publication, presque tout revient au premier niveau. Les cotations de notre teneur de marché se placent derrière
celles des autres (il n'a pas été le premier à atteindre le niveau) et perdent tout de même 57.5 lots-tick par publication en moyenne (\cref{lst:hf:news-and-event-trading:mmloss})~: à 12.50~dollars le
tick, soit 719~dollars.

La seconde moitié du protocole concerne la réinscription. Après la publication, le flux est dense (vingt et une fois le niveau normal au départ, décroissant sur trente secondes) et
toxique (un mark-out de trois ticks au départ, décroissant sur cinq secondes). Coter à partir de $t$ secondes et jusqu'à deux minutes après la publication rapporte
(\cref{fig:hf:news-and-event-trading:reentry})~: 872 lots-tick à partir de la publication elle-même, 911 à partir d'une demi-seconde, 1\,052 à partir de cinq secondes, et moins
au-delà, à mesure que le volume s'éteint. Un teneur de marché qui reste en place pendant la publication gagne $872-58=814$~; celui qui retire ses cotations puis les réinscrit à cinq secondes gagne
1\,052.

\begin{omfigure}
\begin{tikzpicture}
\begin{axis}[width=11cm,height=5.2cm,xlabel={\small délai de réinscription après la publication (secondes)},ylabel={\small lots-tick jusqu'à deux minutes},xmin=0,
  xmax=62,ymin=0,ymax=1150,tick label style={font=\footnotesize},grid=major,grid style={gray!25},table/col sep=comma]
\addplot[omThm,thick,mark=*,mark size=1.3pt] table[x=t,y=value]{figdata/hft/18-news-and-event-trading/reentry.csv};
\draw[omDef,dashed] (axis cs:0,814) -- (axis cs:62,814);
\node[font=\scriptsize,anchor=south east,omDef] at (axis cs:59,814) {rester en place pendant la publication};
\end{axis}
\end{tikzpicture}

\omcaption{Fourchette moins mark-outs gagnée en cotant à partir du moment de réinscription jusqu'à deux minutes après une publication, en lots-tick, lorsque le flux postérieur à la publication
part de vingt et une fois le niveau normal et décroît sur trente secondes, et que son mark-out part de trois ticks et décroît sur cinq~; en tirets, un teneur de marché
qui n'a jamais retiré ses cotations, après sa perte dans la course. Données~: \texttt{hf\_news.reentry\_curve}.}\label{fig:hf:news-and-event-trading:reentry}
\end{omfigure}

\section{Contrôles de risque pour le trading sur événements}

Le trading sur événements concentre le risque en quelques secondes. Les contrôles du chapitre~27 s'appliquent avec des paramètres propres aux événements~: limites de position et de perte qui se resserrent
avant une publication, une vérification que le chiffre analysé se situe dans une fourchette plausible par rapport au consensus avant qu'une stratégie preneuse de liquidité n'agisse dessus, une taille maximale
par titre, et un kill switch qu'un opérateur peut actionner pendant la fenêtre de l'événement. Le protocole du teneur de marché est lui-même un contrôle de risque~: une
annulation automatique à une heure programmée, vérifiée avant l'événement, et non laissée aux réflexes d'un trader.

\section{Fiches de stratégie}

\begin{strategyfile}{Course à la publication sur contrats à terme macroéconomiques}\label{strat:hf:news-and-event-trading:macro}
\sfield{Qui paie, et pourquoi} Les fournisseurs dont les cotations sont périmées pendant les millisecondes qui suivent un chiffre programmé.
\sfield{Instruments et plateformes} Contrats à terme sur taux, sur indices boursiers et sur devises~; le flux officiel de la publication ou celui d'un fournisseur.
\sfield{Signal} La surprise par rapport au consensus, rapportée au mouvement de chaque instrument lors des publications passées.
\sfield{Taille et exécution} Prendre les niveaux périmés situés dans le mouvement attendu, du plus grand au plus petit~; sortir dans le flux postérieur à la publication.
\sfield{Coûts} Flux et colocalisation au point de publication et sur les bourses~; frais~; le risque de signe contraire si la correspondance est fausse.
\sfield{Comment elle meurt} La course~: dans le modèle, le premier niveau capte 88~\% de la valeur~; le deuxième, le plus rapide après lui, n'en obtient que peu.
\sfield{Horizon, capacité, infrastructure} De la microseconde à la seconde~; capacité tirée de la profondeur périmée, que les protocoles réduisent.
\sfield{Backtest honnêtement} L'horodatage exact de la publication chez le trader, la profondeur qui était au repos à ce moment-là, et le concurrent le plus rapide.
\sfield{Sources} Hu, Pan et Wang (2017)~; le présent chapitre.
\end{strategyfile}

\begin{strategyfile}{Course à la publication de résultats lisibles par machine}\label{strat:hf:news-and-event-trading:earnings}
\sfield{Qui paie, et pourquoi} Les fournisseurs dont les cotations sur l'action et sur ses options sont en retard sur la dépêche.
\sfield{Instruments et plateformes} L'action, ses options, son fonds sectoriel~; un flux de dépêches lisibles par machine.
\sfield{Signal} Les champs analysés (résultats, prévisions de la direction) rapportés au consensus.
\sfield{Taille et exécution} De petites tailles, pondérées par la confiance dans l'analyse~; arrêter dès qu'un champ sort de sa plage plausible.
\sfield{Coûts} Le flux~; les erreurs d'analyse (un taux d'erreur de 20~\% transforme des captures de 2~ticks et des pertes de 10~ticks en perte nette).
\sfield{Comment elle meurt} Les erreurs d'analyse et les concurrents plus rapides sur le même flux.
\sfield{Horizon, capacité, infrastructure} De la milliseconde à la minute.
\sfield{Backtest honnêtement} Les horodatages du flux à la réception, et non à la publication~; les dépêches telles qu'elles ont d'abord été envoyées, y compris les corrections.
\sfield{Sources} Le livre~8, chapitre~17.
\end{strategyfile}

\begin{strategyfile}{Réaction aux messages des réseaux sociaux}\label{strat:hf:news-and-event-trading:social}
\sfield{Qui paie, et pourquoi} Les traders plus lents qui réagissent au même message~; ou, sur les faux messages, les traders qui ont agi sur eux.
\sfield{Instruments et plateformes} Les contrats à terme sur indices et les actions citées~; les flux des réseaux sociaux.
\sfield{Signal} Les messages de sources vérifiées, assortis de scores de crédibilité et de la confirmation d'autres sources.
\sfield{Taille et exécution} Très petite tant que l'information n'est pas confirmée~; on parie contre les mouvements non confirmés, comme après le piratage d'avril 2013.
\sfield{Coûts} Les faux positifs~; le flux.
\sfield{Comment elle meurt} Les comptes piratés ou falsifiés~; les actions réglementaires contre les messages manipulateurs.
\sfield{Horizon, capacité, infrastructure} De la seconde à la minute.
\sfield{Backtest honnêtement} Les messages tels qu'ils sont apparus, y compris ceux qui ont été supprimés, avec l'horodatage de réception.
\sfield{Sources} L'encadré daté.
\end{strategyfile}

\begin{strategyfile}{Protocole d'événement du teneur de marché~: retirer, élargir, revenir}\label{strat:hf:news-and-event-trading:protocol}
\sfield{Qui paie, et pourquoi} Le flux important qui suit la publication, une fois sa toxicité retombée.
\sfield{Instruments et plateformes} Chaque instrument coté par le teneur de marché que la publication fait bouger.
\sfield{Signal} Le calendrier des publications~; après la publication, la décroissance des mark-outs.
\sfield{Taille et exécution} Annuler avant la publication~; revenir après les secondes toxiques (cinq dans le modèle) avec un spread plus large, qui se resserre
à mesure que les mark-outs décroissent.
\sfield{Coûts} Le spread abandonné pendant les secondes hors marché.
\sfield{Comment elle meurt} Les événements non programmés, pour lesquels il n'existe pas de calendrier.
\sfield{Horizon, capacité, infrastructure} Les secondes~; un ordonnanceur automatique vérifié avant chaque événement.
\sfield{Backtest honnêtement} Les exécutions et les mark-outs après la publication, seconde par seconde, et non moyennés sur la journée.
\sfield{Sources} Le présent chapitre~: 1\,052 tick-lots revenus à cinq secondes contre 814 restés en place.
\end{strategyfile}

\section{Tutoriel~: huit heures et demie}\label{tut:hf:news-and-event-trading:tut}

\begin{tutorial}
\textbf{Objectif.} Mesurer qui capte les cotations périmées d'une publication, ce que perd un teneur de marché à rester en place, et quand il doit revenir.
\textbf{État final~:} le tableau et la \cref{fig:hf:news-and-event-trading:reentry}.
\begin{tutsteps}
\item \textbf{La course}~: les paliers, dans l'ordre de latence, prennent les niveaux périmés situés dans le mouvement.
\omcode{code/firm/newsrace/firm_newsrace.py}{32}{52}{Le plus rapide d'abord, chaque palier jusqu'à sa capacité~; un lot situé à $k$ ticks de l'ancien prix rapporte $J-k$.}\label{lst:hf:news-and-event-trading:race}
\item \textbf{La perte du teneur de marché} depuis l'arrière de la file de chaque niveau.
\omcode{code/firm/newsrace/firm_newsrace.py}{55}{62}{Nos lots ne sont pris qu'après ceux des autres à chaque niveau.}\label{lst:hf:news-and-event-trading:mmloss}
\item \textbf{Vingt mille publications} avec des surprises normales centrées réduites (\texttt{hf\_news.base}).
\item \textbf{Retour en place}~: le spread diminué des mark-outs décroissants à partir de chaque instant de retour (\texttt{hf\_news.reentry\_curve}).
\end{tutsteps}
\textbf{Pour aller plus loin.} Rejouer la publication sur \texttt{firm.exchsim} avec les fournisseurs comme agents qui réagissent après leurs propres latences~;
ajouter un second instrument que le chiffre fait bouger~; faire estimer la correspondance de la surprise avec une erreur.
\end{tutorial}

\section{Construction~: le module de course à la publication}\label{bld:hf:news-and-event-trading:newsrace}

\begin{build}
\bfield{Rôle} Simuler les publications programmées, la course aux cotations périmées et le \omterm{def:hf:news-and-event-trading:protocol}{protocole d'événement} d'un teneur de marché.
\bfield{Interface} \texttt{move(z, beta)}, \texttt{race(J, stale\_lots, tiers, stale\_ms)},
\texttt{mm\_loss(J, mm\_lots, other\_lots, tiers, stale\_ms)}, \texttt{reentry(t)}, \texttt{simulate(n, seed)}, \texttt{misparse(p\_wrong, capture\_right, loss\_wrong)}. Dépêches non programmées~:
\texttt{firm.newsevent} (livre~8).
\bfield{Règles} Ticks et lots~; les preneurs qui arrivent après la réaction des fournisseurs ne reçoivent rien~; nos lots sont derrière ceux des autres.
\bfield{Tests d'acceptation} \texttt{code/firm/newsrace/tests/}~: une course à 4 ticks calculée à la main, y compris un palier trop lent pour prendre quoi que ce soit~; la perte du
teneur de marché calculée à la main des deux côtés~; le retour après les secondes toxiques l'emporte sur le retour immédiat~; les parts totalisent un et le palier le plus lent n'obtient
rien~; la valeur de l'erreur d'analyse calculée à la main.
\bfield{Pour aller plus loin} Fournisseurs comme agents sur \texttt{firm.exchsim}~; plusieurs instruments~; une courbe de mark-out post-publication apprise.
\end{build}

\begin{omsources}
\item G.~X.~Hu, J.~Pan, J.~Wang, Early peek advantage? Efficient price discovery with tiered information disclosure, \emph{Journal of Financial
Economics} 126(2), 2017, 399--421.
\item New York State Attorney General, communiqué de presse sur Thomson Reuters, 8 juillet 2013.
\item US Securities and Exchange Commission, communiqué de presse 2018-226, 29 septembre 2018.
\item NBC News, reportage sur le piratage du compte Twitter de l'Associated Press, 23 avril 2013.
\end{omsources}

\section{Exercices}

\begin{exercise}[$\star$]\label{exo:hf:news-and-event-trading:1}
Une publication fait bouger le contrat à terme de quatre ticks. Un preneur achète 10 lots à chacun des trois premiers niveaux au-dessus de l'ancien prix. Combien gagne-t-il~?
\end{exercise}

\begin{exercise}[$\star$]\label{exo:hf:news-and-event-trading:2}
Une stratégie fondée sur les dépêches capte 2 ticks quand elle a raison et perd 10 quand elle a tort. À quel taux d'erreur est-elle à l'équilibre~?
\end{exercise}

\begin{exercise}[$\star$]\label{exo:hf:news-and-event-trading:3}
À 12.50~\$ le tick, combien valent 57.5 tick-lots~? Et quel est l'avantage du protocole, de 238 tick-lots par publication, sur douze publications mensuelles~?
\end{exercise}

\begin{exercise}[$\star\star$]\label{exo:hf:news-and-event-trading:4}
Pourquoi le palier le plus rapide prend-il presque tout~?
\end{exercise}

\begin{exercise}[$\star\star$]\label{exo:hf:news-and-event-trading:5}
Pourquoi revenir immédiatement est-il moins bon que revenir après cinq secondes, et pourquoi revenir après trente secondes est-il à nouveau moins bon~?
\end{exercise}

\begin{exercise}[$\star\star$]\label{exo:hf:news-and-event-trading:6}
Hu, Pan et Wang ont montré qu'un accès anticipé pour quelques-uns a pu améliorer la formation des prix. Donner l'argument.
\end{exercise}

\begin{exercise}[$\star\star\star$]\label{exo:hf:news-and-event-trading:7}
\emph{Programmation.} Relancer \texttt{firm.newsrace.simulate} avec des fournisseurs qui mettent à jour après 0.5~ms au lieu de 20~ms. Que deviennent les parts des paliers et
la valeur captée par publication~?
\end{exercise}

\begin{exercise}[$\star\star\star$]\label{exo:hf:news-and-event-trading:8}
\emph{Trouver la faille.} «~Nous avons fait le backtest de notre stratégie de publication sur les horodatages du fournisseur, et elle capte 3 ticks à chaque publication.~»
\end{exercise}

\section{Problème~: Huit heures et demie}

\begin{problem}[{Problème du week-end --- huit heures et demie}]\label{pb:hf:news-and-event-trading:1}
Un teneur de marché sur contrats à terme et plusieurs preneurs de liquidité font face à une publication mensuelle.

\medskip\noindent\textbf{Partie I --- Publications.}
\begin{enumerate}
\item Définir un \omterm{def:hf:news-and-event-trading:lockup}{lock-up de publication} et une \omterm{def:hf:news-and-event-trading:race}{course à la publication}.
\item Que trouvent Hu, Pan et Wang~?
\item Résumer les trois événements de l'encadré daté.
\item Pourquoi le mauvais parsing compte-t-il davantage que la vitesse pour beaucoup de sociétés~?
\end{enumerate}

\medskip\noindent\textbf{Partie II --- La course.}
\begin{enumerate}[resume]
\item Décrire le modèle~: mouvement, niveaux périmés, paliers.
\item Donner les parts des paliers dans la valeur périmée.
\item Pourquoi la perte de notre teneur de marché est-elle calculée depuis le fond de la file~?
\item Que perd-il à une publication, en lots de tick et en dollars~?
\end{enumerate}

\medskip\noindent\textbf{Partie III --- Le protocole.}
\begin{enumerate}[resume]
\item Définir un \omterm{def:hf:news-and-event-trading:protocol}{protocole d'événement}.
\item Donner les valeurs de la courbe de réentrée à 0, 0.5, 5 et 30 secondes.
\item Comparer le maintien des cotations avec leur retrait et une réentrée à cinq secondes.
\item Quels contrôles de risque doivent figurer dans le protocole~?
\end{enumerate}

\medskip\noindent\textbf{Partie IV --- Le verdict.}
\begin{enumerate}[resume]
\item Énoncer le \emph{résultat nommé}~: la part du mouvement de la publication captée par chaque palier de latence, et la perte du teneur de marché avec et sans
son \omterm{def:hf:news-and-event-trading:protocol}{protocole d'événement}.
\item Qui paie la course~?
\item Que changeraient des fournisseurs plus rapides~?
\item Comment le protocole doit-il traiter les événements non programmés~?
\item Pourquoi un backtest doit-il utiliser l'horodatage de réception~?
\item Quelle fiche de stratégie est la plus exposée à une fausse information~?
\item Comment estimeriez-vous la courbe de mark-out après publication~?
\item En une phrase~: que rapporte le \omterm{def:hf:news-and-event-trading:protocol}{protocole d'événement} d'un teneur de marché~?
\end{enumerate}
\end{problem}

\section{Questions d'entretien}

\begin{interviewq}[$\star$ \iqroles{trader}]\label{iq:hf:news-and-event-trading:1}
Les chiffres de l'emploi tombent dans dix minutes. Que faites-vous de vos cotations, et quand revenez-vous~?
\end{interviewq}

\begin{interviewq}[$\star\star$ \iqroles{développeur}]\label{iq:hf:news-and-event-trading:2}
Concevez le chemin qui mène d'une publication lisible par machine à un ordre, et listez les endroits où il peut échouer.
\end{interviewq}

\begin{interviewq}[$\star\star$ \iqroles{chercheur}]\label{iq:hf:news-and-event-trading:3}
Comment associeriez-vous la surprise statistique d'une publication au mouvement de dix instruments différents, et comment testeriez-vous cette association~?
\end{interviewq}

\begin{interviewq}[$\star\star$ \iqroles{risque}]\label{iq:hf:news-and-event-trading:4}
Une stratégie fondée sur les titres d'actualité a gagné de l'argent sur 95~\% des événements et perdu dix fois son gain moyen sur les autres. Comment fixez-vous ses limites~?
\end{interviewq}

\begin{interviewq}[$\star\star$ \iqroles{trader}]\label{iq:hf:news-and-event-trading:5}
Un tweet annonce qu'une société a reçu une offre de rachat. Que vérifiez-vous avant de trader, et dans quel ordre~?
\end{interviewq}

\begin{interviewq}[$\star\star\star$ \iqroles{chercheur}]\label{iq:hf:news-and-event-trading:6}
Avec un flux décroissant en $e^{-s/T_1}$ et des mark-outs en $e^{-s/T_2}$, trouvez le temps de réentrée qui maximise les gains espérés du teneur de marché.
\end{interviewq}
